% Draft sections for the camera-ready. NOT yet \input into submission.tex.
% Numbers marked \pendingnum{...} are filled ONLY from generated artifacts in
% outputs/tables/ (E1/E2/E4), never typed from memory. Working name for the new
% configuration: "calibrated \unilid" (user + co-author decision pending).
% Every number's sentence names its instrument.

% ------------------------------------------------------------------
% (A) Item 3: replacement for the sentence at submission.tex:796.
% Current text: "To avoid zero-probability segmentations when a token's
% expected counts are estimated to be 0 under a given language (most common in
% the few-sample regime), we floor token probabilities at $1e-12$. We did not
% find further smoothing measures to help."
% Replacement:
% ------------------------------------------------------------------

To avoid zero-probability segmentations when a token's expected counts are
estimated to be 0 under a given language (most common in the few-sample
regime), we floor token probabilities at $10^{-12}$ during training. This
training-time floor leaves the resulting probabilities of unseen tokens
miscalibrated across languages: on GlotLID-C, setting them to a shared
constant after training, together with a re-examination step for low-margin
predictions, raises macro F1 substantially (\cref{sec:calibration}).

% ------------------------------------------------------------------
% (B) Item 2: new subsection, proposed placement at the end of
% \section{\unilid} (sec:unilid), after the Computational Complexity paragraph.
% ------------------------------------------------------------------

\subsection{Calibrated \unilid: a shared unseen-token floor and a
re-examination step}\label{sec:calibration}

Two inference-time additions follow from the error analysis of \unilid on
GlotLID-C (\cref{app:perf-breakdown}). We call the resulting configuration
\defn{calibrated \unilid}. Training is unchanged; both additions operate on
the learned per-language distributions and on quantities computed from each
language's own training data.

\paragraph{A shared floor for unseen tokens.} After training, the probability
that language $\lang$ assigns to a token unseen in its training corpus is a
byproduct of the training-time floor and renormalization, and it varies by
language: the median across the 1,940 GlotLID-C languages is $e^{-17.66}$,
with smaller training corpora producing higher unseen-token probabilities.
Because \cref{eq:lang-id} compares scores across languages, these differing
values act as per-language offsets on every unseen token: strings built of
tokens a language has never seen are assigned to it more often when its
unseen-token probability is higher. Calibrated \unilid sets the
log-probability of every unseen token
to a single shared constant, $-21$ (natural log), in every language. Languages
whose training-derived value was higher are penalized more for unseen tokens;
languages whose value was lower are penalized less. On the held-out portion of
the GlotLID-C test pool defined in \cref{app:protocol}, this one change moves
macro F1 from \pendingnum{0.9117-source-paper-eval} to
\pendingnum{floor21-judge-source-paper-eval}.

\paragraph{Re-examining low-margin predictions into under-resourced
languages.} False positives concentrate in languages with little training
data: their distributions are close to their initialization, so they assign
middling probability to many strings, and the shared floor above does not
remove this effect for the smallest corpora. Calibrated \unilid therefore
re-examines a prediction when (i) the predicted language has fewer than
18{,}000 training samples, or is one of four large-corpus languages whose
token distributions are unusually flat for their script (Scots, Banjar,
Aragonese, and West Flemish), and (ii) the top score exceeds the runner-up
score by less than a per-language threshold. The threshold for language
$\lang$ is a percentile of the score margins that $\lang$ achieves on its own
training samples, so it requires no data from other languages: the percentile
is $q_\lang = 5\,(1 - N_\lang/18{,}000)$ for a language with $N_\lang$
training samples, and a fixed 5th percentile for the four flat languages. A
re-examined prediction is reassigned to the highest-ranked of the candidates
ranked 2 to 5 that has at least 100{,}000 training samples and a score within
21 natural-log units of the top candidate; if no candidate qualifies, the
prediction is unchanged.

\paragraph{Provenance of the constants.} The floor $-21$ was chosen from a
sweep over $\{-17, -19, -21, -23\}$ on a held-out validation subset of
GlotLID-C; the proximity bound $21$ from a grid over $0.5$ to $100$ in steps
of $1$ (the optimum is a plateau spanning roughly $15$ to $35$) on a
development subset of the test pool; the $100{,}000$-sample bar from an
analysis of where reassignments send their false positives. Membership of the
four flat languages uses one diagnostic computed on validation data. All
constants are frozen before evaluation on the held-out reporting subset, which
was excluded from every one of these selections; \cref{app:protocol} gives the
full
development protocol and an accounting of which data selected which constant.

\paragraph{Effect.} On GlotLID-C, calibrated \unilid raises macro F1 from
\pendingnum{baseline-fullpool} to \pendingnum{calibrated-fullpool} (full test
set) while lowering macro FPR from \pendingnum{baseline-fpr} to
\pendingnum{calibrated-fpr}; on the held-out subset that no constant was
selected on, the paired bootstrap interval of the improvement over \unilid is
\pendingnum{ci-vs-baseline} (\cref{app:protocol}). The gain is concentrated
in the languages where \cref{tab:resource-tier} reports \unilid's weakness:
for languages with under 1{,}000 training samples, mean per-language F1 rises
from \pendingnum{tail-baseline} to \pendingnum{tail-calibrated} (held-out
subset, all false positives counted). Both additions preserve the
incremental-addition property of \cref{sec:unilid}: adding a language requires
its own training data, the shared constants, and one flatness diagnostic, but
no retraining of other languages. Two scoping notes. First, the gains are
measured on natural-distribution test data, where false-positive volume
matters; on balanced, content-controlled evaluation sets these gains shrink or
reverse, and \cref{app:protocol} reports those measurements. Second, the
re-examination step targets under-resourced languages, so on benchmarks
restricted to larger languages the added mechanisms rarely fire.

% ------------------------------------------------------------------
% (C) Item 4 sample: one appendix paragraph (development protocol opening),
% shown as the style sample for the appendix material.
% ------------------------------------------------------------------

\section{Development protocol for calibrated \unilid}\label{app:protocol}

The constants of \cref{sec:calibration} were selected on data disjoint from
the subset used to report their effect. From the 45{,}377{,}279 scored lines
of the GlotLID-C test pool, we removed two balanced draws (188{,}061 and
185{,}204 lines, at most 100 lines per language, seeds 101 and 201) and split
the remaining 45{,}004{,}014 lines by a seeded 40/60 partition (seed 301)
into a development part (18{,}001{,}573 lines) and a held-out part
(27{,}002{,}441 lines). The floor value was selected on a retired 250{,}000-line
validation sample outside this pool; the proximity bound on the development
part; the re-examination thresholds on each language's own training samples.
The held-out part was used only to confirm each candidate configuration once,
and it is the subset on which the paired bootstrap intervals in this appendix
are computed. Table 1 reports the full test pool for comparability with the
other systems; the table below repeats the comparison on the held-out part.
